\documentclass[conference,a4paper]{IEEEtran}
\IEEEoverridecommandlockouts

\usepackage{cite}
\usepackage{amsmath,amssymb,amsfonts}
\usepackage{graphicx}
\usepackage{xcolor}
\usepackage[hidelinks]{hyperref}
\usepackage{microtype}

\begin{document}

\title{IEEE P3335, Standard for Architecture and Interfaces for Time Card:\\ Motivation and Progress Report}

\author{\IEEEauthorblockN{Ahmad Byagowi}
\IEEEauthorblockA{Time Appliance Project\\
Chair, IEEE P3335 Time Card Working Group\\
byagowi@turbalance.com}
% Add working-group co-authors as needed
\thanks{This paper describes the unapproved IEEE P3335 ; content
is subject to change.}
}

\maketitle

\begin{abstract}
IEEE P3335 defines the generic architecture and interfaces of a Time Card:
a modular timing subsystem that provides traceable time-of-day to a host
system and, via protocols such as PTP and NTP, to networked distributed
systems. The project generalizes the Open Compute Project (OCP) Time Card,
the core of the Open Time Server, into a vendor-neutral standard. We report
the motivation for the work and its progress as of mid-2026: a structurally
complete, pre-ballot working draft covering conformance, architecture,
performance declaration, timing and control interfaces, environment, and
application guidance. The presentation summarizes the draft's key design
decisions (a unified-timescale invariant, stackable conformance profiles,
a transport-neutral control model, and declaration-based performance) and
invites input from the ISPCS community during the working group's 
consideration of the draft.
\end{abstract}

\begin{IEEEkeywords}
IEEE P3335, Time Card, IEEE 1588, PTP hardware clock, time synchronization,
conformance, figures of merit
\end{IEEEkeywords}

\section{Motivation}
Data centers, telecommunications, industrial control, finance, and
scientific measurement increasingly depend on precise, traceable
time-of-day, yet a typical host lacks an adequate local clock, clock
steering, and time-transfer machinery. The OCP Time
Card~\cite{ocp-timecard} (a PCIe card combining a GNSS receiver, a
high-stability holdover oscillator, and an FPGA) showed that an open,
modular timing subsystem can anchor an Open Time
Server~\cite{byagowi2022timecard}, and a diverse ecosystem of commercial
and open implementations followed. Each uses product-specific
signal definitions, control models, and performance claims: integration
effort is duplicated for every card--host pairing, and independent
comparison of products is impractical.

IEEE P3335~\cite{p3335par} (PAR approved December 2022, sponsored by the
IEEE Instrumentation and Measurement Society) addresses this gap with three
goals: plug-and-play interoperability between Time Card implementations and
the hosts and operating systems that support them, including POSIX
integration via a shared PTP hardware clock (PHC); modular implementation
around defined building blocks (time source, local oscillator, time
processor); and unambiguous comparison through common figures of merit,
measurement points, and operating conditions. The working group's guiding principle is
to embrace implementation diversity while ensuring interoperability.

\section{Architecture}
The draft defines a Time Card by externally observable behavior rather than
physical construction: one conformance model covers a discrete add-in card,
embedded gateway, an FPGA or SoC function, and an external timing unit. Internally, the architecture names a local timing function (the oscillator
state, from quartz to miniaturized atomic references), a time generator, a
timestamping facility, and receive, providing, and control interfaces
(Fig.~\ref{fig:arch}); the timing data plane is separated from management
and control.

\begin{figure}[t]  % Figure is Illegible
\centering
\includegraphics[width=0.85\columnwidth]{figures/architecture-context.pdf}
\caption{Time Card context: external references enter through receive
interfaces, discipline one unified timescale, and reach the host through
providing and control interfaces.}
\label{fig:arch}
\end{figure}

Recognized receive interfaces include GNSS (multi-constellation, with
antenna-state reporting), PTP per IEEE 1588-2019~\cite{ieee1588} or IEEE
802.1AS-2025~\cite{ieee8021as}, NTP~\cite{rfc5905}, IRIG AM/DCLS, discrete
1PPS and frequency inputs, and other time-transfer methods such as White
Rabbit and wireless time services. Receive-less designs (free-running or
holdover-only) are permitted when declared, and multi-reference designs
must document their selection policy (priority, validity, failure
detection, fallback, recovery).

The central invariant is the unified timescale: a conforming instance
generates exactly one, and every providing interface (1PPS, frequency
outputs, time-of-day messages, PTP, NTP, IRIG, PCIe Precision Time
Measurement (PTM), and host-bus timestamps) publishes, encodes, or derives
from it, with a declared alignment error for each time- or phase-bearing
interface. Ensemble operation must collapse to one selected or synthesized
reference for disciplining that timescale. Behavior at
leap events, holdover entry and exit, reference failover, and manual phase steps
must be documented, and continuous-phase outputs may not take undocumented
phase steps while locked.

Host time transfer receives explicit treatment: mappings must document
discovery, driver-visible registers or APIs, timestamp formats, and the
latency, asymmetry, and correction data needed to evaluate time transfer
into the host. PCIe PTM implementations must
state which measurement point PTM timestamps represent and the correction
terms relating PTM time to other providing interfaces, addressing the
historically underspecified last segment between the timing card and the
host operating system.

\section{Conformance and Performance Declaration}
Base conformance requires at least one timing providing interface, one
unified timescale, one accessible control interface, and full documentation
of interfaces, performance metrics per measurement point, environmental
limits, and validity conditions (warm-up, lock time, reference,
temperature, measurement bandwidth). Conditional profiles
stack on top: a Physical Timing Output profile with concrete 1PPS limits
(SMA connector into a 50~$\Omega$ load, rising-edge on-time mark,
10\%--90\% rise time below 10~ns, alignment error below 100~ns between the
1PPS output and interfaces claiming alignment to it), a PCIe host-mapping
profile, and Managed and Secure Infrastructure security profiles. The
Managed profile requires authenticated, role-separated management with
integrity-checked updates; Secure Infrastructure adds signed update
payloads verified against a trust anchor, rollback protection, secure or
measured boot before the timing service enters the locked state, and
hardware-backed key storage. Claiming any optional feature binds the
implementation to that feature's requirements.

Instead of universal numeric product classes, the draft mandates bounded,
reproducible performance declarations. Each claim states its metric,
measurement point, operating mode (locked, holdover, free-running, warm-up,
failover, recovery), reference, environmental conditions, observation
interval, uncertainty, and decision rule, with methods per IEEE
1139-2022~\cite{ieee1139}, IEEE 1193-2022~\cite{ieee1193}, and ITU-T
G.810/G.8260~\cite{itug8260}, and results metrologically traceable to a
national metrology institute or a declared reference. MTIE is required for
bounded time-error and holdover claims; TDEV, ADEV, phase noise, and a
defined pulse-jitter statistic apply where relevant. Metrics may not be
used interchangeably, and typical values cannot replace bounded ones.
Informative annexes add formula-level metric definitions, a holdover error
model, and example test procedures.

\section{Control Interfaces}
A transport-neutral baseline information model keeps object semantics
identical whether control rides SMBus, I2C, or I3C, PCIe MMIO, IPMI or
NC-SI, or REST, gRPC, Redfish, or SNMP. Required objects cover identity and
versioning, device state (initializing through locked, holdover, degraded,
and fault), timescale identity, sync-source status, alarms, security state,
and coherent (tear-free) time reads. Conditional telemetry spans phase error in picoseconds, fractional
frequency offset in parts in $10^{15}$, MTIE estimates, holdover elapsed
time, and GNSS status. Configuration writes must report acceptance and whether they
can cause a phase step, loss of lock, or output interruption.

\section{Progress, Process, and Remaining Work}
The working group formed from the OCP time-appliances community after PAR
approval; contributions on architecture, host integration, use cases, and
metrics accumulated through 2025, and drafting moved
to a public git repository~\cite{p3335repo} in October 2025. As of July
2026 the draft is structurally complete and pre-ballot: Clauses 1--10 plus
metrics and test-procedure annexes compile to a 48-page document carrying
411 machine-indexed requirement statements. The draft is written as per-clause Markdown with
pull-request-based comment resolution and a generated requirements index
for traceability; the comment round from the April 2026 working-group
meeting is merged, and all 14 normative references passed a technical use
audit. Open pre-ballot decisions include ratifying the 1PPS electrical and
alignment limits, choosing between a mandated cross-vendor host ABI (a
common register layout or OS API) and transport-neutral semantics alone,
freezing the baseline control vocabulary, ratifying normative-reference
editions, fixing the security-profile cryptographic suites, and an
independent timing-metrology review. The presentation will detail these decisions, share experience running an
IEEE standards project on open git tooling, and solicit review and
participation from the ISPCS community while the draft is considered in the working group.

\bibliographystyle{IEEEtran}
\bibliography{references}

\end{document}
